\section{Active-tail localization, information, and original-policy bounds}
\label{supp:activetail119}

\subsection{Proof of Theorem~\ref{thm:tailinformation119}}
\begin{proof}
Because $f$ is continuous, the fundamental theorem of calculus gives
$H'(u)=f(u+L)-f(u)=D(u)$.  For convex $f$, its increment over a fixed
positive length is nondecreasing.  One direct proof is to use the
almost-everywhere derivative of a Lipschitz convex function:
\[
 D(v)-D(u)=\int_u^v\{f'(s+L)-f'(s)\}\,ds\ge0
 \quad(u<v).
\]
The derivative is monotone wherever defined, and changing its values on
a null set does not change the integral.  Also,
\[
 |D(v)-D(u)|\le |f(v+L)-f(u+L)|+|f(v)-f(u)|
                  \le2K|v-u|.
\]
Thus $H$ is convex with a $2K$-Lipschitz derivative.  Compactness
gives a minimizer on $I$, and the derivative signs give the asserted
endpoint rules.

We prove that a minimizing central interval discards precisely a lower
$\alpha$ fraction of the outcomes, allowing equality at the threshold.
At an interior minimizer $u^*$, $f(u^*)=f(u^*+L)=\eta$.  Convexity
gives $f\le\eta$ between these endpoints and $f\ge\eta$ outside.
The outside set has original probability $t$.  The feasible density
$q=t^{-1}\boldsymbol1_{[-R,u^*]\cup[u^*+L,R]}$ therefore has
\[
 \int fq\,dP=\eta+t^{-1}\mathbb E(f-\eta)_+.
\]
The equality is unaffected by any flat portion at level $\eta$.
For an arbitrary feasible AVaR density $\widetilde q$, the constraints
$0\le\widetilde q\le t^{-1}$ and $\int\widetilde q\,dP=1$ give
\[
 \int f\widetilde q\,dP
 =\eta+\int(f-\eta)\widetilde q\,dP
 \le\eta+t^{-1}\mathbb E(f-\eta)_+.
\]
Consequently this selected $q$ is optimal.

At the left endpoint, $f(-R+L)\ge f(-R)$.  Convexity implies that
the maximum on $[-R,-R+L]$ is $f(-R+L)$ and that $f$ is
nondecreasing to the right of $-R+L$; hence the same threshold
argument applies with $\eta=f(-R+L)$.  At the right endpoint, the
inequality is reversed and its reflected argument applies.  These
arguments include equal endpoint values and constant functions.
They establish \eqref{eq:exacttail119} without a density for the
distribution of $f(Z)$ or a no-atom assumption.

When neither endpoint rule holds, continuity of $D$ gives an interior
zero.  Sign bisection maintains a bracket containing at least one zero,
whose width after $m$ halvings is $2Rt\,2^{-m}$ because $|I|=2Rt$.
Let $u^*$ be a zero in that bracket.  Integrating the $2K$-Lipschitz
derivative from $u^*$ to its midpoint $u_m$ gives
\[
 0\le H(u_m)-H(u^*)
 \le K|u_m-u^*|^2
 \le KR^2t^2\,4^{-m}.
\]
The total integral of $f$ is constant in $u$, so
$\operatorname{AVaR}_\alpha(f)-R_{u_m}(f)
 =[H(u_m)-H(u^*)]/(2Rt)$.
This proves \eqref{eq:taillocalization119}.  An endpoint or exact-zero
exit has zero loss, so no smoothness-boundary term is omitted.

For a nonempty cell $[l,r]$ of length $h$ and midpoint $z$, convexity
gives the pointwise secant upper bound and Jensen's lower bound,
\[
 h f(z)\le\int_l^r f(s)\,ds
       \le\tfrac h2\{f(l)+f(r)\}.
\]
Their difference is
\[
 \frac{h^2}{4}\left\{
       \frac{f(r)-f(z)}{r-z}-\frac{f(z)-f(l)}{z-l}\right\}.
\]
The braces are nonnegative.  When all cells are ordered from left
to right, convexity orders every successive secant slope, including
across the omitted central interval.  The sum of the bracketed slope
differences is therefore at most the difference between the largest
and smallest of these slopes, hence at most $2K$.  With maximal
cell length $h_{\max}$, summation and normalization by $2Rt$ give
\[
 \overline R-\underline R
 \le\frac{K h_{\max}^2}{4Rt}
 \le\frac{KRt}{n^2}.
\]
This proves the interval and its width bound.

If the selected tail lengths are $b_1,b_2\ge0$, then
$b_1+b_2=2Rt$.  Use $\lceil nb_j/(2Rt)\rceil$ cells in each
nonempty tail.  Since the two unrounded counts sum to the integer
$n$, their ceilings sum to at most $n+1$.  Midpoints and distinct
endpoints require at most $2(n+1)+2=2n+4$ values, allowing two
nonempty components.  The two endpoint sign checks require at most
four values, and $m$ bisections require at most $2m$ more.  Reusing
duplicates can only reduce this count.  The query bound follows.
Putting $m=\lceil\log_2n\rceil$ gives $4^{-m}\le n^{-2}$.

For completeness, the lower bound preserves its original deterministic
adaptive quantifiers.  Let $q=-R+2R\alpha$ and $b=R-q=2Rt$,
and run the algorithm on
\[
 f_0(z)=\frac{K}{2b}(z-q)_+^2.
\]
This function is convex, nondecreasing, and $K$-Lipschitz.  Insert
the $N$ queried locations lying strictly between $q$ and $R$, together
with the two endpoints.  Their distinct consecutive interval lengths
are $h_1,\ldots,h_s$, where $s\le N+1$ and $\sum_i h_i=b$.
Define $f_1$ to equal $f_0$ to the left of $q$ and its chord on each
of these intervals.  The chord slopes are nondecreasing in $[0,K]$;
thus $f_1$ belongs to the same class and agrees with every query
answer.  Induction on the adaptive transcript shows that the algorithm
cannot distinguish $f_0$ and $f_1$.

Since both functions are nondecreasing, their AVaRs are their averages
over $[q,R]$.  The integral of the chord error on an interval of
length $h_i$ is $K h_i^3/(12b)$.  Jensen's inequality gives
\[
 \operatorname{AVaR}_\alpha(f_1)-\operatorname{AVaR}_\alpha(f_0)
 =\frac{K}{12b^2}\sum_i h_i^3
 \ge\frac{Kb}{12(N+1)^2}
 =\frac{KRt}{6(N+1)^2}.
\]
A sound common interval contains both values.  This proves the lower
bound.  With $\alpha=0$, the same construction uses $q=-R$ and
the upper algorithm simply partitions the full interval; no central
interval or tail sign search is necessary.  For example, for every
$N\ge32$ one may take $n=\lfloor N/4\rfloor$ in the upper
algorithm; its stated query bound is at most $N$ and its width is at
most $(75/2)KRt/N^2$.  Together with the lower bound this proves the
claimed uniform order, without concealing a logarithmic factor in the
error rate.
\end{proof}

\subsection{Primitive curvature and inexact tail localization}
\begin{proof}
Write $\mu_x=(4-5\beta/4)/d$.  The inherited primitive curvature
bound says that the jointly convex last-date objective satisfies the
strong-convexity inequality with terms
$\mu_x\|x-y\|_2^2/2$ and $|a-b|^2$.
For any $x,y$, select their optimal actions $a_x,a_y$.
The capacity $A$ is affine, so
$(1-\theta)a_x+\theta a_y$ is feasible at
$(1-\theta)x+\theta y$ for $0\le\theta\le1$.
Evaluate the strong-convexity inequality at this feasible interpolated
action and then take the value infimum.  Dropping the nonnegative
action-distance term proves
\[
 V_1((1-\theta)x+\theta y)
 \le(1-\theta)V_1(x)+\theta V_1(y)
       -\tfrac12\mu_x\theta(1-\theta)\|x-y\|_2^2.
\]
Thus $V_1$ is $\mu_x$-strongly convex.  For fixed original $x,a$,
the successor map is affine in $z$ with direction
$s=(1,-1,\ldots,1,-1)$ and $\|s\|_2^2=d$.
Composition yields $m_z=\mu_xd=181/64$.  The inherited state
modulus $15/d$ in the $\ell^1$ norm gives
$|f_a(z)-f_a(z')|\le15|z-z'|$.

For $u<v$, strong convexity of $f_a$ gives monotonicity of
$f_a'(z)-m_zz$ at its almost-everywhere derivative points.  Therefore
\[
 D_a(v)-D_a(u)
 =\int_u^v\{f_a'(s+L)-f_a'(s)\}\,ds
 \ge m_zL(v-u).
\]
This is precisely $m_zL$-strong convexity of $H_a$.  For the true
$d=D_a(u)$ it follows that
\[
 H_a(v)\ge H_a(u)+d(v-u)+\tfrac12m_zL(v-u)^2
 \quad(v\in I).
\]
Taking the infimum proves that the normalized tail deficit is at most
the expression in \eqref{eq:taildefect119} at the true $d$.
The negative of an infimum of affine functions of $d$ is convex;
its maximum over $[d_-,d_+]$ is bounded by the larger endpoint
value.  Completing the square gives the projected minimizer.  These
facts prove \eqref{eq:taildefect119} for outward derivative intervals,
including both constrained endpoints and $d=0$.
\end{proof}

\subsection{Proof of Theorem~\ref{thm:investmenttail119}}
\begin{proof}
The selected density is exactly feasible under the original law.
Integrating the midpoint affine supports over their symmetric cells
cancels the term proportional to $z-z_i$, yielding
$R_u(f_a)\ge r_-$.  Convexity of $f_a$ and its endpoint upper
values give $R_u(f_a)\le r_+$.  By the tail-defect bound,
\[
 r_-\le\operatorname{AVaR}_\alpha(f_a)
                 \le r_++\kappa(u).
\]
The complete terminal policy has
$V_1\le J_1\le V_1+\varepsilon_1$ throughout the original state
domain.  Monotonicity and cash invariance, for the same original
risk primitive, imply
\[
 \operatorname{AVaR}_\alpha(f_a)
 \le\operatorname{AVaR}_\alpha(J_1\circ F_a)
 \le\operatorname{AVaR}_\alpha(f_a)+\varepsilon_1.
\]
Adding unchanged stage cost and discounting proves
\eqref{eq:tailpolicyinterval119}.  No convexity of $J_1$ has been used.

For another original feasible action $v$, the transition satisfies
$F(x,v,z)=F(x,a,z)+b(v-a)$.  Therefore each global support gives
\[
 V_1(F(x,v,z))
 \ge\ell_i+p_i^{\mathsf T}\{b(v-a)+s(z-z_i)\}
 \quad(z\in B_i).
\]
Integrating over the \emph{same fixed} feasible density removes the
last term.  The original stage cost satisfies
$h(v)\ge h(a)+h'(a)(v-a)+(v-a)^2$ because $h''\ge2$.
Combining these inequalities yields
\[
 Q(v)\ge W+s(v-a)+(v-a)^2
 \quad(0\le v\le A(x)).
\]
Taking the infimum gives \eqref{eq:tailoptimum119}; the true
$V_0$ is this original optimal-future infimum.  For uncertain slope,
the constrained quadratic minimum is concave as a function of its
slope, so its minimum over an interval is bounded below by the
smaller endpoint minimum.  Taking outward endpoints preserves every
direction of inequality.  A smaller exact feasible action domain can
be enlarged when constructing a lower bound, but may never be shrunk
to obtain a favorable minimum.

Since $V_0\le J(a)$, the policy upper minus the optimum lower is
a valid nonnegative regret allowance.  The same inference holds at
$a=0$ and $a=A(x)$: no first-order condition for an interior actor
has been assumed.  Tail search need not succeed, a fitted critic
need not equal $V_1$, and a numerical candidate optimizer need not
return an optimum, provided the displayed outward witnesses remain
valid and the realized action remains feasible.
\end{proof}

\subsection{Paid precision and conditional allocation}
\begin{proof}
Qualification at a chosen stage of a finite sequence follows from its
own verified inequalities; no monotonicity of earlier reported widths
is needed.  Summing the actual cost of every attempted stage is the
complete cost.  If a deterministic ladder has work $cn_j$ per stage
with $n_j=2^jn_0$, its total work through $n_J$ is less than
$2cn_J$.  Such a bound requires this work model and does not remove
fixed costs, descendant solves, refused trials, or rational fallback.

For the allocation formula write $e_b=B4^{-m}$ and
$e_q=A/n^2$.  At a nonboundary continuous optimum the positive
constraint binds, hence $e_q=\epsilon-e_b$ and the objective is
\[
 F(e_b)=c_q\sqrt{\frac A{\epsilon-e_b}}
             +\frac{c_b}{\log4}\log\frac B{e_b}.
\]
Its second derivative is positive for $0<e_b<\epsilon$.
Its derivative tends to negative infinity at zero and positive
infinity at $\epsilon$.  It therefore has exactly one minimizer,
and setting its derivative equal to zero gives
\eqref{eq:tailallocation119}.  If this point violates $n\ge1$ or
$m\ge0$, the convex problem in $(n,m)$ is minimized on the
corresponding boundary or at the remaining feasible stationary point;
the constraints $A/n^2+B4^{-m}\le\epsilon$ form a convex set.
Increasing either resolution decreases its allowance, so rounding
both upward preserves feasibility and adds less than $c_q+c_b$.
All coefficients must be fixed properties of the declared graph and
cost model.  Estimating them after observing a favorable outcome does
not turn the resulting allocation into a prospective experiment.
\end{proof}

\subsection{Recursive interpretation}
Suppose at each date the constructed certificate supplies a uniform
original-optimal-continuation action gap $\gamma_t\ge0$ for a
feasible deployed policy.  If its next-date policy loss is at most
$e_{t+1}$, the unchanged Bellman preference gives
\[
 J_t^\pi\le Q_t(V_{t+1}^*;x,\pi_t(x))+\beta_te_{t+1}
          \le V_t^*+\gamma_t+\beta_te_{t+1}.
\]
Backward induction from the declared terminal discrepancy gives the
dated discounted sum.  A support of an NBO's own fitted continuation
instead requires the retained one-sided or reference-centered
continuation-error account.  It cannot be silently promoted to a
support of $V_{t+1}^*$.  The currently executed original two-date
construction obtains the needed $V_1$ supports from the primitive
last-date oracle, and retains its actual terminal-controller allowance.
